\subsection{Implementation Details}
\label{sec:impl}

We implement \sys as a kernel backend for PyTorch. A PyTorch program can be compiled into a \sys mega-kernel via PyTorch’s compilation interface by specifying \sys as the backend, i.e., {\tt torch.compile(backend=\sys)}. This call invokes the \sys compiler, which generates a mega-kernel and returns it as a callable PyTorch function. Executing this function issues a single launch of the generated mega-kernel. The current \sys implementation consists of approximately 40K lines of C++, 84K lines of CUDA, and 10K lines of Python. The in-kernel parallel runtime is written in CUDA and uses semaphores in device memory to coordinate workers and schedulers. The \sys compiler, implemented in C++ and Python, automatically transforms an input tensor program into an optimized \graph tailored to specific GPU types. For compute tasks, the compiler integrates the Mirage superoptimizer~\cite{wu2025mirage} to automatically generate optimized CUDA implementations and uses NVSHMEM~\cite{nvshmem} to support in-kernel inter-GPU communication.

%We implement \sys as a kernel backend for PyTorch. A PyTorch program can be directly compiled into a \sys mega-kernel through PyTorch's compilation interface by specifying \sys as the kernel backend: {\tt torch.compile(backend=\sys)}. This invocation triggers the \Sys compiler to generate a mega-kernel and returns the mega-kernel to users, who can subsequently call this mega-kernel. An invocation to this mega-kernel will execute the mega-kernel. \Sys is implemented in xxx lines of code in C++, CUDA, and Python. The in-kernel parallel runtime system was implemented in CUDA, using semophere in device memory to enable synchronization across workers and schedulers. The \Sys compiler is implemented in C++ and Python to automatically transform an input tensor program into highly optimized \graph for execution on the parallel runtime. The compiler also leverages the Mirage superoptimizer to generate CUDA implementation for all compute tasks in the \graph~\cite{wu2025mirage}, and directly uses NVSHMEM~\cite{nvshmem} APIs for in-kernel inter-GPU communications.

Our implementation also includes several key optimizations to minimize runtime overhead and efficiently support dynamic workloads.

\paragraph{Task-launch overhead.}
Because \sys decomposes computation into tasks that are substantially finer-grained than traditional GPU kernels, minimizing per-task overhead is essential for performance. \sys employs several techniques to keep task-launch costs low.
First, the runtime uses extremely lightweight workers and schedulers: event and task queues are implemented as circular buffers in GPU device memory, and enqueue/dequeue operations rely solely on low-cost {\tt atomicAdd} instructions.
Second, \sys adopts a decentralized scheduling strategy in which each scheduler assigns tasks using only local state. This design avoids global coordination and eliminates the communication and synchronization overheads inherent to globally coordinated scheduling.

While \sys currently uses decentralized scheduling, the runtime is designed to support alternative strategies, including globally coordinated scheduling, with only minor code modifications. Exploring these designs and their performance trade-offs is an interesting direction for future work.

%\paragraph{Task overhead.} \Sys's tasks involve much finer granularity than kernels and as a result it's critical to minimize task launch overhead in \Sys. \Sys involves several impelemntation optimizations to minimize task launch overhead. First, the \Sys runtime uses extremely lightweight workers and schedulers. Both event queues and task queues are implemented using circular buffers on GPU device memory, and enqueue and dequeue operations are implemented using the {\tt atomicAdd} operation. Second, our current implementation uses {\em decentralized scheduling} where schedulers assign tasks to workers using only local information. This approach avoids expensive inter-scheduler communications and synchronizations. 

\paragraph{Supporting runtime dynamism.}
To demonstrate \sys’s ability to support highly dynamic workloads, we extend the system with mechanisms required for LLM serving, including {\em continuous batching}~\cite{yu2022orca} and {\em paged attention}~\cite{kwon2023vllm}. When processing the start event of a task graph, the scheduler prepares a new decoding iteration by (1) removing completed requests from the previous iteration, (2) admitting newly arrived requests, and (3) updating per-request KV-cache metadata. All of this logic executes within a single task, and the KV-cache metadata is stored in device memory for direct use by paged-attention tasks.

To handle the dynamic batch sizes intrinsic to LLM serving, \sys generates multiple \graphs specialized for representative batch sizes (powers of two up to the maximum batch). At runtime, the scheduler selects the appropriate graph based on the current batch size. This approach allows the compiler to produce \graphs optimized for specific batch sizes while still preserving flexibility for highly dynamic workloads.

%\paragraph{Support runtime dynamism.} To demonstrate \sys's ability to support dynamic workloads, we use LLM serving as the main use case and have enhanced \sys to support {\em continuous batching}~\cite{yu2022orca} and {\em paged attention}~\cite{kwon2023vllm}, two critical optimizations for LLM serving. To support these optimizations, the scheduler processing the start event of a task graph prepares a new serving iteration, including (1) dropping finished requests from the previous batch, (2) adding newly arrived requests for the next batch, and (3) updating KV cache pages for active requests. All processing was performed in a single task and KV cache page metadata was stored in a tensor in device memory, which will be used by all paged-attention tasks. 

%To support dynamic batch sizes, \sys generates multiple task graphs for representative batch sizes (i.e., all powers of two until reaching the maximum batch size). Based on the dynamic batch size, the scheduler will launch the corresponding task graph. This design allows the \sys compiler to generate \graphs optimized for specific batch sizes. 

%\begin{itemize}
%    \item Support vllm and pytorch frontend (\ZJ{TODO: @Jianan})
%    \item Support both online and offline llm serving
%\end{itemize}
